\chapter{Complex Sentences and Clause Chaining}\label{ch:complex}

\section{Overview: clause-combining strategies}

Iridian employs several strategies for combining clauses. The dominant strategy is \textbf{converb chaining}, in which all non-final clauses are expressed as non-finite converb forms linked to a single main (finite) clause. Additional strategies include coordination with conjunctions, quotative constructions for reported speech, and \ird{še}-clauses for a limited set of complement and adjunct types.

The centrality of the converb to Iridian clause-combining means that what would be expressed as multiple finite clauses in a Slavic language typically appears in Iridian as a chain of converbs followed by a single finite verb. This has consequences for information structure (only the final, finite verb carries aspect and alignment), for switch-reference tracking, and for discourse organization (see Chapter~\ref{ch:discourse}).

\section{Converb chains: structure and alignment}\label{sec:converb-chains}

\subsection{The finite-verb-final rule}

In a converb chain, all verbs except the last are non-finite (converbs). Only the final verb is finite and carries aspect, mood, and alignment. This rule is absolute in standard written Iridian; in colloquial and dialectal speech, finite verbs may occasionally appear in non-final positions in certain coordinated structures.

\subsection{Accusative alignment within converb clauses}

Case marking within converb clauses always follows nominative-accusative alignment: the absolutive argument is the agent or sole argument, and the patient appears in the indirect (oblique) case. The ergative case is never found inside a converb clause, because converbs carry no aspect and the ergative alignment is triggered only by the perfective and retrospective aspects.

\subsection{Switch-reference and topic continuity}\label{sec:switch-reference}

By default, the subject of a converb clause is coreferent with the subject of the following clause (same-subject, SS reading). When the subjects differ (different-subject, DS), a pronominal clitic is attached to the converb to mark the switch. If both clauses have full nominal subjects, the full NPs disambiguate and no clitic is required.

\section{Sequential constructions}\label{sec:sequential}

\subsection{Forms: \ird{-u} and \ird{-ení}}

The sequential converb has two forms with subtly different temporal force. The \ird{-u} form implies a temporally bounded, completed sub-event preceding the main clause event. The \ird{-ení} form is temporally underspecified and is compatible with both sequential and simultaneous readings.

\pex
\a
\begingl
\gla Marek zdravení Janek znovime.//
\glb Marek sleep-\Seq{} Janek study-\Prog{}//
\glft \trsl{While Marek sleeps, Janek is studying.} \lingcontext{simultaneous reading with \ird{-ení}}//
\endgl
\xe

\subsection{The copula converb and the discourse particle \ird{nu}}\label{sec:seq-copula}

The marked copular verb \ird{neh-} (see §\,\ref{sec:neh} of Chapter~\ref{ch:simple}) has the sequential converb form \ird{nesu}, a suppletive form established independently of the regular aspect paradigm. In colloquial speech, \ird{nesu} is contracted to \ird{nu}. The form \ird{nu} has been grammaticalized in some dialects as a discourse particle inserted between clauses in the manner of English `and then', without strong sequential implication and without retaining the inchoative/copular force of the underlying verb.

\subsection{Interaction with applicative and causative}

\pex
\a \ird{piašt-u} / \ird{piašt-ení} \trsl{having eaten / and [then] ate\ldots}
\a \ird{piašt-éb-u} / \ird{piašt-éb-ní} \trsl{having fed someone}
\a \ird{piašt-nau} / \ird{piašt-na-ní} \trsl{having made someone eat}
\xe

\section{Temporal constructions}\label{sec:temporal}

\subsection{The simultaneous / manner converb (\ird{-ěc})}

The simultaneous converb (\ird{-ěc}) describes an action that is ongoing or co-occurring with the main clause event, and also functions as a manner converb. The suffix belongs to the palatalizing domain of front-vocalic morphology (see Section~\ref{sec:front-palatalization}), not to the special sibilant-replacement system of the antipassive. When the converb and main clause have different subjects, the sequential \ird{-ení} is preferred for purely simultaneous meaning.

\subsection{The six temporal converbs}

\begin{table}[ht]
	\footnotesize\sffamily
	\caption{Temporal converb suffixes.}\label{tab:tempconverbs}
	\medskip
	\begin{tabularx}{0.9\textwidth}{L l L}
		\toprule
		{\sc suffix} & {\sc meaning} & {\sc function} \\
		\midrule
		\ird{-ezak}    & until              & main clause holds until converb-clause event \\
		\ird{-eškady}  & when, around when  & loose temporal co-occurrence \\
		\ird{-eeby}    & since              & from the time of the converb-clause event \\
		\ird{-ešhoume} & as soon as         & immediately upon the converb-clause event \\
		\ird{-eebymy}  & after              & after the converb-clause event \\
		\ird{-esny}    & before             & before the converb-clause event \\
		\bottomrule
	\end{tabularx}
\end{table}

\section{Causal constructions}\label{sec:causal}

% ============================================================
% STATUS: STUB — causal converb suffix not yet documented.
% The 2026 notes list converb types 1, 2, 4, 5 and skip 3 (causal).
% This section is a placeholder until the suffix is confirmed.
% ============================================================

A causal converb introduces the reason or cause for the main clause event (`because X, Y'). The suffix is not yet confirmed in the current revision and will be documented here once established.

\section{Concessive constructions}\label{sec:concessive}

The concessive converb (\ird{-ouc}) introduces a conceded or contrasted clause (`although X, Y'). It licenses the irrealis mood in the main clause when the conceded content is contrary to expectations.

\section{Conditional constructions}\label{sec:conditional}

\subsection{The four conditional converbs}

\begin{table}[ht]
	\footnotesize\sffamily
	\caption{Conditional converb suffixes.}\label{tab:conditional}
	\medskip
	\begin{tabularx}{0.8\textwidth}{L l L}
		\toprule
		{\sc suffix} & {\sc meaning} & {\sc notes} \\
		\midrule
		\ird{-ic}   & if (possibility)       & does not trigger palatalization \\
		\ird{-emy}  & if not (possibility)   & \\
		\ird{-ezmy} & if/when (expectation)  & \\
		\ird{-ena}  & if/when not (expected) & \\
		\bottomrule
	\end{tabularx}
\end{table}

\subsection{No palatalization with \ird{-ic}}\label{sec:cond-nopala}

Despite beginning with the front vowel \ird{-i-}, the conditional converb \ird{-ic} does not trigger palatalization of the thematic consonant. This is a historical accident: the suffix was grammaticalized during a period when this \ird{-i-} lacked palatalizing force. The conditional is therefore distinguishable from the irrealis (\ird{-ěv}), the simultaneous converb (\ird{-ěc}), and the conjunctive linker \ird{-e}, all of which belong to the productive palatalizing domain.

\section{Purposive constructions}\label{sec:purposive}

The purposive converb (\ird{-it}) expresses purpose (`in order to X') and serves as the complement of modal and desiderative verbs. Palatalization is irregular, applying mainly to Class~I stems.

\subsection{The intentional future reanalysis}

In colloquial speech, the purposive converb has been reanalyzed as a future/intentional marker with first-person clitics when no governing verb is present:

\ex
\begingl
\gla Dá trave piašcitem.//
\glb \textsc{1sg} bread.\Ind{} eat-\Purp{}-\textsc{1sg}//
\glft \trsl{I'm gonna eat bread.} \lingcontext{colloquial intentional future}//
\endgl
\xe

\section{Contrastive constructions}\label{sec:contrastive}

The contrastive converb (\ird{-ema}) introduces a contrasting clause (`X, but/whereas Y'), covering simple contrast, adversative contrast, and parallel contrast.

\section{Coordination}\label{sec:coordination}

% ============================================================
% STATUS: STUB — systematic treatment needed.
% Key conjunctions: a (additive), še (additive / complementizer),
% má / ozná (adversative), disjunctives.
% ============================================================

Coordination is expressed by \ird{a} (additive), \ird{še} (additive, also complementizer), \ird{má} / \ird{ozná} (adversative), and a set of disjunctives.

\section{Quotative constructions and reported speech}\label{sec:quotative}

% ============================================================
% STATUS: STUB — quotative complementizer to-že, direct quotatives,
% evidential/quotative verbal suffix. Port from 2020 archive.
% ============================================================

Reported speech uses the quotative complementizer \ird{to-že} for indirect quotation and allows fronted direct quotations before a \textit{verbum dicendi}. Clause-linking and certain quotative constructions further employ the verbal linker \ird{-e}, whose morphophonemics are described in Section~\ref{sec:conjunctive-linker} of Chapter~\ref{ch:verbs}. The evidential/quotative morpheme on the verb signals attribution to another speaker.

\section{Epistemic particles and evidential extensions}\label{sec:evidential-particles}

% ============================================================
% STATUS: STUB — see open question in ROADMAP.md §3.
% ============================================================

\section{Fixed expressions with converbs}\label{sec:fixed-converbs}

% ============================================================
% Thanks, apologies, congratulations, greetings that are historically
% grammaticalized converb forms — to be documented.
% ============================================================
